\begin{proposition}[A secondary Hopf surface with torsion-free first homology]\hypertarget{tag-01F9}{}\label{cdp:st:C.3.3}
\StatementRecord{cdp:st:C.3.3}{proved}{01F9}{the definitions and constructions stated in the proof}


There is a secondary Hopf surface $Q_H$ with
$H_1(Q_H;\Z)=\Z$ and $H^2(Q_H;\Z)=0$.
\end{proposition}


\begin{proof}

\begin{enumerate}[leftmargin=*]
\item Write a quaternion as $(q_0,q_1,q_2,q_3)$ and put
$\varphi=(1+\sqrt5)/2$. Let $G$ be the following set of 120 unit
quaternions: the eight signed coordinate units; the sixteen points
$(\pm1,\pm1,\pm1,\pm1)/2$; and the 96 points obtained by even
coordinate permutations of
$(0,\pm1,\pm\varphi,\pm\varphi^{-1})/2$, with independent signs.
Quaternion multiplication makes this the binary icosahedral group.
Multiplication of the listed coordinates is closed in $\Q(\sqrt5)^4$,
and the commutator map on this finite group is surjective.
Every element is therefore a commutator, and $G^{\rm ab}=0$.

\item Use the faithful holomorphic representation
\[
(a,b,c,d)\longmapsto
\begin{pmatrix}a+ib&c+id\\-c+id&a-ib\end{pmatrix}\in SU(2)
\]
and let the generator of $\Z$ act on $\C^2\setminus\{0\}$ by
$v\mapsto2v$. The action of $\Z\times G$ is free. An equality
$2^ngv=v$ first gives $n=0$ by norms, and a unit quaternion whose
matrix has eigenvalue one is the identity. Proper discontinuity
follows from the radial scaling and finiteness of $G$.
Polar coordinates identify the quotient smoothly with
\[
Q_H=(\C^2\setminus\{0\})/(2^{\Z}\times G)
\simeq (\mathbb R/(\log2)\Z)\times(S^3/G).
\]
It is compact, and its universal cover is $\C^2\setminus\{0\}$.
The fundamental group is $\Z\times G$, whose abelianization is
$\Z$. The quotient by $2^{\Z}$ alone is a primary Hopf surface
covering $Q_H$ with degree 120. The fundamental group of $Q_H$ is
noncyclic, so $Q_H$ is secondary Hopf. The product description gives
$e(Q_H)=0$, and the same duality and universal-coefficient argument
as before gives $H^2(Q_H;\Z)=0$.

\end{enumerate}
\end{proof}